\documentclass{situationist-preprint}

\title{Censure and Hard Treatment}
\msid{ETH-26-b}
\msfamily{ETH-26}
\msdiscipline{Normative and Applied Ethics}
\mstype{Companion manuscript}
\mskeywords{punishment; censure; hard treatment; communication; penance; expressivism}

\begin{document}
\maketitle

\begin{abstract}
Communicative theories of punishment face the question why censure must be conveyed by hard treatment rather than by words. Three answers are on offer: hard treatment as a conventional vehicle of condemnation, as a prudential supplement to moral reasons for compliance, and as a penance. We argue, developing the reply of ETH-26-a, that hard treatment is the medium in which censure expresses the seriousness of a wrong, and that this is why it is required, bounded above, and absent when no one is wronged. The conventional answer leaves the medium arbitrary and the prudential answer abandons the communicative justification; the account given here explains both what is right in the penance view and the limits of hard treatment.
\end{abstract}

\section{Three accounts}

\citet{Feinberg1965} held that punishment expresses condemnation and that hard treatment is the conventional vehicle of that expression in our society. On this account another vehicle would serve as well if convention changed. \citet{vonHirsch1993} holds that censure is conveyed by the penal sanction, but that hard treatment is justified also as a prudential supplement, giving those insufficiently moved by moral reasons a further reason to comply. \citet{Duff2001} holds that hard treatment is a secular penance through which the offender can come to understand and repent her wrong and express that repentance to others.

\section{The argument}

ETH-26-a defines a punishment as justified if it communicates censure of the offence and the offender has reason to accept the censure (its Definition~2.4). Its Assumption~2.6 requires that a punishment communicating censure of an offence $o$ bring cumulative hard treatment to the seriousness $w(o)$ at some period; its Assumption~2.5 requires that a punishment whose censure the offender has reason to accept never exceed $k\cdot w(o)$.

\begin{proposition}\label{p:medium}
Hard treatment is required for the communication of censure because the seriousness of a wrong is expressed by the burden it is answered with, and a message that varies with the seriousness of wrongs must vary in a dimension that offender, victim and community read as burden.
\end{proposition}

\begin{proof}[Argument]
Censure that represents a wrong accurately represents its seriousness, since a censure that treated a grave wrong as minor would misrepresent it. Words can state that a wrong is grave, but a statement costs its addressee nothing, and the same words can be uttered of any wrong. What shows that a wrong is taken seriously is what the response asks of the one who committed it: the time, liberty or resources she must give up. This is why Assumption~2.6 of ETH-26-a ties censure to cumulative burden, and why the burden must scale with $w(o)$.
\end{proof}

Proposition~\ref{p:medium} keeps what is right in the conventional account, that hard treatment expresses rather than merely causes, while denying that the medium is arbitrary: burdens are not a conventional sign of seriousness but its natural measure. It rejects the prudential supplement, since a reason to comply that does not depend on the offender's accepting the censure is not part of the communicative justification, and a punishment set by its deterrent effect would have no upper bound drawn from the offence. It agrees with the penance account that hard treatment gives the offender something to do with the censure, and adds a constraint that account needs: a penance that the offender has reason to accept is bounded by the seriousness of the wrong. Hence Theorem~3.3(a) of ETH-26-a, and hence its termination result.

\begin{corollary}\label{c:words}
Where an act wrongs no one, censure has nothing to express by burden, and justified punishment imposes none.
\end{corollary}

\begin{proof}
This is Corollary~3.5 of ETH-26-a, and Proposition~\ref{p:medium} explains it: if $w(o)=0$, the burden that represents the seriousness of $o$ is zero, and by Assumption~2.5 any burden exceeds what the offender has reason to accept.
\end{proof}

\begin{objection}
If hard treatment is the medium of censure, the unrepentant offender, who does not read it as censure, is subjected to burden that communicates nothing.
\end{objection}

\begin{reply}
Definition~2.4 of ETH-26-a requires that the offender have reason to accept the censure, not that she accept it. The burden communicates the censure whether or not it is heard, as a sincere apology addressed to someone who ignores it is still addressed to her \citep{Bennett2008}. What the unrepentant offender lacks is uptake, and the bound on hard treatment does not depend on uptake.
\end{reply}

\section{Conclusion}

Hard treatment is neither a conventional sign of censure nor a prudential addition to it. It is the dimension in which censure measures a wrong, and its role as measure is what both requires it and limits it.

\begin{thebibliography}{9}\small
\bibitem[Bennett(2008)]{Bennett2008} Bennett, C. (2008). \emph{The Apology Ritual: A Philosophical Theory of Punishment}. Cambridge University Press.
\bibitem[Duff(2001)]{Duff2001} Duff, R.~A. (2001). \emph{Punishment, Communication, and Community}. Oxford University Press.
\bibitem[Feinberg(1965)]{Feinberg1965} Feinberg, J. (1965). The expressive function of punishment. \emph{The Monist}, 49(3), 397--423.
\bibitem[von Hirsch(1993)]{vonHirsch1993} von Hirsch, A. (1993). \emph{Censure and Sanctions}. Oxford University Press.
\end{thebibliography}

\end{document}
